% ============================================================
%  style/griffbilder.tex — Griffbilder (Akkorddiagramme) definieren und zeichnen
%  (wird von style.tex eingebunden)
% ============================================================
% ============================================================
%  Griffbilder (Akkorddiagramme)
% ============================================================
\usepackage{xparse}
\usepackage{etoolbox}

\newlength{\griffsw}\setlength{\griffsw}{2.0mm} % Saitenabstand
\newlength{\griffsh}\setlength{\griffsh}{2.5mm} % Bundabstand
\newcommand{\griffnamefont}{\serif\bfseries\fontsize{\SzGriffName}{\LdGriffName}\selectfont\color{inkheading}}
\newcommand{\grifffingerfont}{\serif\bfseries\fontsize{6.5}{6.5}\selectfont\color{inkmain}}
\newcommand{\griffmarkfont}{\serif\bfseries\fontsize{7}{7}\selectfont\color{inkchords}}
% Leeres Griffbrett zum handschriftlichen Eintragen eigener Akkorde
\newcommand{\leergriff}{
  \griff{}{.,.,.,.,.,.}{}%
}

\def\GriffStop{}
\def\griffbarreparse#1:#2:#3\GriffStop{%
  \draw[inkchords,line width=4.2pt,line cap=round]
  (#2-1,#1-0.42) -- (#3-1,#1-0.42);
}
\def\griffdotparse#1:#2:#3\GriffStop{%
  \fill[inkchords] (#1-1,#2-0.5) circle (2.3pt);
  \notblank{#3}{\node at (#1-1,4.35) {\grifffingerfont #3};}{}%
}

% #1 Name  #2 Topmarken (6x, kommagetrennt: x/o/.)
% #3 Punkte (Saite:Bund:Finger, kommagetrennt, Finger optional)
% #4 Startbund-Label (optional), #5 Barré Bund:StartSaite:EndSaite (optional)
%
% \griff     zeichnet das Griffbild (und merkt es sich beim Einlesen von chords.tex)
% \defgriff  merkt sich ein Griffbild nur für den NÄCHSTEN Song (z.B. eine
%            besondere Version wie bei Wonderwall). VOR \begin{song} in die
%            Song-Datei schreiben; nach \end{song} gilt wieder der normale Griff.
% \griffzeichnen  eigentliche Zeichenroutine
%
% Griffe für \griffe kommen aus chords.tex (wird am Ende als Tabelle gedruckt)
% und griffe-zusatz.tex (wird nur eingelesen, nicht gedruckt).
\newif\ifgriffsammeln
\NewDocumentCommand{\griffbauen}{ m m m o o }{%
  \IfNoValueTF{#4}{\def\grifftmp{\griffzeichnen{#1}{#2}{#3}}}{%
    \IfNoValueTF{#5}{\def\grifftmp{\griffzeichnen{#1}{#2}{#3}[#4]}}%
                    {\def\grifftmp{\griffzeichnen{#1}{#2}{#3}[#4][#5]}}}%
}
\newcommand{\griffspeichern}[1]{%
  \ifblank{#1}{}{\global\cslet{griffdef:\detokenize{#1}}\grifftmp}%
}
\NewDocumentCommand{\griff}{ m m m o o }{%
  \IfNoValueTF{#4}{\griffbauen{#1}{#2}{#3}}{%
    \IfNoValueTF{#5}{\griffbauen{#1}{#2}{#3}[#4]}{\griffbauen{#1}{#2}{#3}[#4][#5]}}%
  \ifgriffsammeln\griffspeichern{#1}\else\grifftmp\fi
}
\NewDocumentCommand{\defgriff}{ m m m o o }{%
  \IfNoValueTF{#4}{\griffbauen{#1}{#2}{#3}}{%
    \IfNoValueTF{#5}{\griffbauen{#1}{#2}{#3}[#4]}{\griffbauen{#1}{#2}{#3}[#4][#5]}}%
  \edef\grifflokalname{\detokenize{#1}}%
  \xifinlist{\grifflokalname}{\griffelokal}{}{%
    \ifcsdef{griffdef:\grifflokalname}%
      {\global\csletcs{griffbak:\grifflokalname}{griffdef:\grifflokalname}}{}%
    \listxadd{\griffelokal}{\grifflokalname}}%
  \griffspeichern{#1}%
}
% songspezifische Griffe am Songende wieder auf den Normalgriff zurücksetzen
\newcommand{\griffelokal}{}
\newcommand{\grifflokalreset}[1]{%
  \ifcsdef{griffbak:#1}%
    {\global\csletcs{griffdef:#1}{griffbak:#1}%
     \global\expandafter\let\csname griffbak:#1\endcsname\undefined}%
    {\global\expandafter\let\csname griffdef:#1\endcsname\undefined}%
}
\newcommand{\griffelokalzuruecksetzen}{%
  \forlistloop{\grifflokalreset}{\griffelokal}\gdef\griffelokal{}%
}

% Beim Dokumentstart alle Griffe aus chords.tex und griffe-zusatz.tex
% einsammeln (ohne Ausgabe), damit sie in jedem Song per Name abrufbar sind.
% (entfällt bei \griffeaustrue)
\AtBeginDocument{%
  \ifgriffeaus\else
  \begingroup
    \griffsammelntrue
    \let\\\relax \let\hfill\relax
    \setbox0=\vbox{\hsize=\textwidth % ============================================================
%  Griffbilder (Saiten-Zuordnung korrigiert)
% ============================================================

% C-Familie
\griff{C}{x,.,.,o,.,o}{2:3:3,3:2:2,5:1:1} \hfill
\griff{Cm}{x,.,.,.,.,.}{4:3:3,3:3:4,5:2:2}[3][1:2:6] \hfill
\griff{C7}{x,.,.,.,.,o}{2:3:3,3:2:2,4:3:4,5:1:1} \hfill
\griff{Cm7}{x,.,.,.,.,.}{3:3:4,5:2:2}[3][1:2:6] \hfill
\griff{Cmaj7}{x,.,.,o,o,o}{2:3:3,3:2:2} \hfill
\griff{C/B}{x,.,.,o,.,o}{2:2:3,3:2:2,5:1:1} \\[0mm]

% C#- / Db-Familie
\griff{C\#}{x,.,.,.,.,.}{4:3:3,3:3:4,5:3:2}[4][1:2:6] \hfill
\griff{C\#m}{x,.,.,.,.,.}{4:3:3,3:3:4,5:2:2}[4][1:2:6] \hfill
\griff{C\#7}{x,x,.,.,.,.}{3:3:3,4:4:4,5:2:1,6:4:4} \hfill
\griff{C\#m7}{x,.,.,.,.,.}{3:3:4,5:2:2}[4][1:2:6] \hfill
\griff{C\#sus4}{x,x,.,.,.,.}{3:4:4,4:1:1,5:2:2,6:2:3} \hfill
\griff{C\#dim7}{x,x,.,.,.,.}{3:2:2,4:3:3,5:2:1,6:3:4} \\[0mm]

% D-Familie
\griff{D}{x,x,o,.,.,.}{4:2:1,5:3:3,6:2:2} \hfill
\griff{Dm}{x,x,o,.,.,.}{4:2:2,5:3:3,6:1:1} \hfill
\griff{D7}{x,x,o,.,.,.}{4:2:2,5:1:1,6:2:3} \hfill
\griff{Dmaj7}{x,x,o,.,.,.}{4:2:1,5:2:2,6:2:3} \hfill
\griff{Dsus4}{x,x,o,.,.,.}{4:2:1,5:3:3,6:3:4} \hfill
\griff{Dsus2}{x,x,o,.,.,o}{4:1:1,5:3:3} \\[0mm]

% D#- / Eb-Familie
\griff{D/F\#}{.,x,.,.,.,x}{1:2:1,4:2:2,5:3:3} \hfill
\griff{D\#m}{x,x,.,.,.,.}{3:1:1,4:3:3,5:4:4,6:2:2} \hfill
\griff{D\#7}{x,x,.,.,.,.}{4:3:2,5:2:1,6:3:3} \hfill
\griff{D\#maj7}{x,x,.,.,.,.}{3:1:1,4:3:3,5:3:4,6:3:2} \hfill
\griff{D\#sus4}{x,x,.,.,.,.}{4:3:2,5:4:4,6:4:3} \hfill
\griff{D\#dim7}{x,x,.,.,.,.}{3:1:1,4:2:2,5:1:1,6:2:3} \\[0mm]

% E-Familie
\griff{E}{o,.,.,.,o,o}{2:2:2,3:2:3,4:1:1} \hfill
\griff{Em}{o,.,.,o,o,o}{2:2:2,3:2:3} \hfill
\griff{E7}{o,.,o,.,o,o}{2:2:2,4:1:1} \hfill
\griff{Em7}{o,.,.,o,.,.}{2:2:1,3:2:2,5:3:3,6:3:4} \hfill
\griff{Esus4}{o,.,.,.,o,o}{2:2:2,3:2:3,4:2:4} \hfill
\griff{Edim7}{x,x,.,.,.,.}{3:2:2,4:3:3,5:2:1,6:3:4} \\[0mm]

% F-Familie
\griff{F}{.,.,.,.,.,.}{3:3:3,2:3:4,4:2:2}[][1:1:6] \hfill
\griff{Fm}{.,.,.,.,.,.}{3:3:3,2:3:4}[][1:1:6] \hfill
\griff{F7}{.,.,.,.,.,.}{2:3:4,4:2:2}[][1:1:6] \hfill
\griff{Fmaj7}{x,x,.,.,.,o}{3:3:3,4:2:2,5:1:1} \hfill
\griff{Fsus4}{x,x,.,.,.,.}{3:3:3,4:3:4,5:1:1,6:1:1} \hfill
\griff{Fdim7}{x,x,o,.,o,.}{4:1:1,6:1:2} \\[0mm]

% F#- / Gb-Familie
\griff{F\#}{.,.,.,.,.,.}{3:3:3,2:3:4,4:2:2}[2][1:1:6] \hfill
\griff{F\#m}{.,.,.,.,.,.}{3:3:3,2:3:4}[2][1:1:6] \hfill
\griff{F\#7}{.,.,.,.,.,.}{2:3:4,4:2:2}[2][1:1:6] \hfill
\griff{F\#maj7}{x,x,.,.,.,.}{3:4:4,4:3:3,5:2:2,6:1:1} \hfill
\griff{F\#sus4}{x,x,.,.,.,.}{3:4:4,4:4:4,5:2:2,6:2:2} \hfill
\griff{F\#dim7}{x,x,.,.,.,.}{3:1:1,4:2:2,5:1:1,6:2:3} \\[0mm]

% G-Familie
\griff{G}{.,.,o,o,.,.}{5:3:3,1:3:3,2:2:2,6:3:4} \hfill
\griff{Gm}{.,.,.,.,.,.}{3:3:3,2:3:4}[3][1:1:6] \hfill
\griff{G7}{.,.,o,o,o,.}{1:3:3,2:2:2,6:1:1} \hfill
\griff{Gmaj7}{.,.,o,o,o,.}{1:3:3,2:2:2,6:2:1} \hfill
\griff{Gsus4}{.,.,o,o,o,.}{1:3:3,2:2:2,6:3:4,5:1:1} \hfill
\griff{G/B}{x,.,o,o,o,.}{2:2:2,6:3:4}\\[0mm]

% G#- / Ab-Familie
\griff{G\#}{.,.,.,.,.,.}{3:3:3,2:3:4,4:2:2}[4][1:1:6] \hfill
\griff{G\#m}{.,.,.,.,.,.}{3:3:3,2:3:4}[4][1:1:6] \hfill
\griff{G\#7}{x,x,.,.,.,.}{2:3:4,4:2:2}[4][1:1:6] \hfill
\griff{G\#maj7}{x,x,.,.,.,.}{3:1:1,4:1:1,5:1:1,6:3:3}[4] \hfill
\griff{G\#sus4}{x,x,.,.,.,.}{3:1:1,4:1:1,5:2:2,6:4:4}[4] \hfill
\griff{G\#dim7}{x,x,o,.,o,.}{4:1:1,6:1:2} \\[0mm]

% A-Familie
\griff{A}{x,o,.,.,.,o}{3:2:1,4:2:2,5:2:3} \hfill
\griff{Am}{x,o,.,.,.,o}{3:2:2,4:2:3,5:1:1} \hfill
\griff{A7}{x,o,.,o,.,o}{3:2:2,5:2:3} \hfill
\griff{Amaj7}{x,o,.,.,.,o}{3:2:2,4:1:1,5:2:3} \hfill
\griff{Asus4}{x,o,.,.,.,o}{3:2:1,4:2:2,5:3:3} \hfill
\griff{Adim7}{x,x,.,.,.,.}{3:1:1,4:2:2,5:1:1,6:2:3} \\[0mm]

% A#- / Bb-Familie
\griff{A\#}{x,x,.,.,.,.}{3:3:3,4:3:4,5:3:2,6:1:1} \hfill
\griff{A\#m}{x,x,.,.,.,.}{3:3:3,4:3:4,5:2:2,6:1:1} \hfill
\griff{A\#7}{x,x,.,.,.,.}{3:3:3,4:1:1,5:3:4,6:1:1} \hfill
\griff{A\#maj7}{x,x,.,.,.,.}{3:3:3,4:2:2,5:3:4,6:1:1} \hfill
\griff{A\#sus4}{x,x,.,.,.,.}{3:3:2,4:3:3,5:4:4,6:1:1} \hfill
\griff{A\#dim7}{x,x,.,.,.,.}{3:2:2,4:3:3,5:2:1,6:3:4} \\[0mm]

% B- / H-Familie
\griff{B}{x,.,.,.,.,.}{3:3:4,4:3:3,5:3:2}[2][1:2:6] \hfill
\griff{Bm}{x,.,.,.,.,.}{3:3:4,4:3:3,5:2:2}[2][1:2:6] \hfill
\griff{B7}{x,.,.,.,.,.}{2:2:2,3:1:1,4:2:3,6:2:4} \hfill
\griff{Bm7}{x,.,o,.,.,.}{2:2:2,4:2:3,6:2:4} \hfill
\griff{Bmaj7}{x,x,.,.,.,.}{3:4:4,4:3:2,5:4:3,6:2:1} \hfill
\griff{Bsus4}{x,x,.,.,.,.}{3:4:3,4:4:4,5:5:5,6:2:1} \\[0mm]% ============================================================
%  Zusätzliche Griffbilder für \griffe{...} unter den Songs
%  Gleiche Syntax wie chords.tex, wird aber NICHT als Tabelle gedruckt.
%  #2 Topmarken tiefe E -> hohe e, #3 Saite:Bund:Finger (Saite 1 = tiefe E)
% ============================================================

% Erweiterte Akkorde
\griff{Cadd9}{x,.,.,o,.,.}{2:3:2,3:2:1,5:3:3,6:3:4}          % x32033
\griff{Csus4}{x,.,.,o,.,.}{2:3:3,3:3:4,5:1:1,6:1:1}          % x33011
\griff{Dm7}{x,x,o,.,.,.}{4:2:2,5:1:1,6:1:1}                  % xx0211
\griff{E7sus4}{o,.,o,.,o,o}{2:2:2,4:2:3}                     % 020200
\griff{F6}{x,x,.,.,.,.}{3:3:3,4:2:2,5:3:4,6:1:1}             % xx3231
\griff{Fadd9}{x,x,.,.,.,.}{3:3:3,4:2:2,5:1:1,6:3:4}          % xx3213
\griff{Fm6}{x,x,o,.,.,.}{4:1:1,5:1:1,6:1:1}                  % xx0111
\griff{F\#m7}{.,.,.,.,.,.}{2:3:3}[2][1:1:6]                  % 242222
\griff{Bsus2}{x,.,.,.,.,.}{3:3:3,4:3:4}[2][1:2:6]            % x24422
\griff{Bb}{x,x,.,.,.,.}{3:3:3,4:3:4,5:3:2,6:1:1}             % xx3331 (= A#)

% Slash-Akkorde (Basston)
\griff{C/E}{o,.,.,o,.,o}{2:3:3,3:2:2,5:1:1}                  % 032010
\griff{C/G}{.,.,.,o,.,o}{1:3:3,2:3:4,3:2:2,5:1:1}            % 332010
\griff{Dm/F}{x,x,.,.,.,.}{3:3:3,4:2:2,5:3:4,6:1:1}           % xx3231
\griff{D7/F\#}{.,x,o,.,.,.}{1:2:,4:2:2,5:1:1,6:2:3}          % 2x0212 (Daumen)
\griff{E7/G\#}{.,x,o,.,o,o}{1:4:3,4:1:1}                     % 4x0100
\griff{F/A}{x,o,.,.,.,.}{3:3:3,4:2:2,5:1:1,6:1:1}            % x03211
\griff{F/C}{x,.,.,.,.,.}{2:3:3,3:3:4,4:2:2,5:1:1,6:1:1}      % x33211
\griff{F\#m/E}{o,.,.,.,.,.}{2:3:3,3:3:4}[2][1:2:6]           % 044222
\griff{G/D}{x,x,o,o,o,.}{6:3:3}                              % xx0003
\griff{G/F\#}{.,x,o,o,.,.}{1:2:1,5:3:3,6:3:4}                % 2x0033
\griff{Bm/A}{x,o,.,.,.,.}{3:3:3,4:3:4,5:2:2,6:1:1}[2]        % x04432

% Powerchords
\griff{E5}{o,.,.,x,x,x}{2:2:2,3:2:3}                         % 022xxx
\griff{F5}{.,.,.,x,x,x}{1:1:1,2:3:3,3:3:4}                   % 133xxx
\griff{F\#5}{.,.,.,x,x,x}{1:1:1,2:3:3,3:3:4}[2]              % 244xxx
\griff{G5}{.,.,.,x,x,x}{1:1:1,2:3:3,3:3:4}[3]                % 355xxx
\griff{G\#5}{.,.,.,x,x,x}{1:1:1,2:3:3,3:3:4}[4]              % 466xxx
}%
  \endgroup
  \fi
}

\NewDocumentCommand{\griffzeichnen}{ m m m o o }{%
  \begin{tikzpicture}[x=\griffsw,y=-\griffsh,baseline=(current bounding box.south)]
    \node[anchor=south,align=center] at (2.5,-1.05) {\griffnamefont #1};
    \foreach \i in {0,...,5}{
        \draw[rulecolor,thin] (\i,0) -- (\i,4);
      }
    \IfNoValueTF{#4}{%
      \draw[inkmain,line width=1pt] (0,0) -- (5,0);
    }{%
      \draw[rulecolor,thin] (0,0) -- (5,0);
      \node[anchor=east] at (-0.3,0.5) {\griffmarkfont #4};
    }%
    \foreach \j in {1,...,4}{
        \draw[rulecolor,thin] (0,\j) -- (5,\j);
      }
    \def\marklist{#2}
    \foreach \m [count=\i from 0] in \marklist{
      \ifdefstring{\m}{x}{\node at (\i,-0.32) {\griffmarkfont $\times$};}{}%
      \ifdefstring{\m}{o}{\draw[inkchords,thin] (\i,-0.32) circle (2pt);}{}%
    }
    \IfNoValueF{#5}{%
      \griffbarreparse#5\GriffStop
    }%
        \ifblank{#3}{\node at (2.5,4.35) {\grifffingerfont\phantom{1}};}{%
      \def\dotlist{#3}%
      \foreach \d in \dotlist{
        \expandafter\griffdotparse\d\GriffStop
      }
    }%
  \end{tikzpicture}%
}

